% Generated draft: requires mathematical and source review.
\documentclass[11pt]{article}
\usepackage[margin=25mm]{geometry}
\usepackage{fontspec}
\setmainfont{Latin Modern Roman}
\usepackage{amsmath,amssymb,mathtools,unicode-math}
\setmathfont{Latin Modern Math}
\usepackage{enumitem,needspace,xurl,hyperref}
\setlength{\parindent}{0pt}
\setlength{\parskip}{5pt}
\setlength{\emergencystretch}{3em}
\newcommand{\N}{\mathbb{N}}
\newcommand{\Z}{\mathbb{Z}}
\newcommand{\Q}{\mathbb{Q}}
\newcommand{\R}{\mathbb{R}}
\newcommand{\C}{\mathbb{C}}
\newcommand{\F}{\mathbb{F}}
\newcommand{\A}{\mathbb{A}}
\newcommand{\PP}{\mathbb{P}}
\newcommand{\E}{\mathbb{E}}
\newcommand{\eps}{\varepsilon}
\newcommand{\dd}{\,\mathrm{d}}
\begin{document}
\section*{UnsolvedMath ID: 30002090}
\textbf{Source problem code:} catalogue entry
\subsection{Definitions and mathematical statement}
In the usual diagrammatic setting, let \(\mathcal D\) be a collection of knot diagrams equipped with Reidemeister moves, and let \(C(D)\) denote the crossings of \(D\in\mathcal D\). A binary parity assigns a map
\[
p_D:C(D)\longrightarrow \Z/2\Z
\]
to each diagram. Corresponding crossings retain their labels under moves; a crossing removed by a first Reidemeister move has label \(0\); the two crossings removed by a second Reidemeister move have equal labels; and the three crossings participating in a third Reidemeister move have labels whose sum is \(0\). These definitions describe the usual binary framework for interpreting the question; the supplied record does not specify the intended group-valued extension or its axioms.

A knot in a \(3\)-manifold \(M\) is an embedding \(S^1\hookrightarrow M\), with the category and equivalence relation to be specified. For a link with a fixed finite collection of circle components, a link homotopy is a homotopy of its component maps during which images of distinct components remain disjoint, while self-intersections of an individual component are permitted.

\textbf{Research questions.} Develop parity theory in the following directions:
\begin{enumerate}
\item Describe how group information can be assigned to crossings and how a corresponding group bracket can be defined. The source does not specify the group data, the compatibility rules under diagram moves, the bracket's target, or the identities the bracket must satisfy.
\item Determine whether an appropriate parity theory exists for knots in arbitrary \(3\)-manifolds. The intended scope ranges over arbitrary \(3\)-manifolds \(M\), but the source supplies no restrictions concerning orientability, boundary, compactness, or category, and supplies no diagram model or required properties of the theory. In particular, in any diagrammatic setting admitting the binary definition above, the assignment \(p_D(c)=0\) is a parity; the source does not state what nontriviality or additional information is required.
\item Find applications of parity theory to link homotopy theory. The source does not identify a particular ambient manifold, invariant, classification problem, or other required application.
\end{enumerate}
These are open-ended research directions, not a fully specified existence theorem. Their unresolved choices must be fixed before a unique precise mathematical proposition can be extracted.
\subsection{Short English statement}
The problem asks how to encode group information at crossings and define a group bracket, extend parity theory to knots in arbitrary \(3\)-manifolds, and apply it to link homotopy. The supplied record labels these directions open, but leaves the required axioms, nontriviality, and intended applications unspecified.
\subsection{Sources}
\begin{itemize}
\item Supplied local UnsolvedMath registry record, permanent ID \(30002090\), title \emph{Group-Valued Parity and Knots in Three-Manifolds}, exact source code \texttt{catalogue entry}. The record lists proposed year \(2012\), gives no individual proposer, and labels the statement \texttt{exact} and its status \texttt{open}. Its statement-verification field reports checking the mathematical content against the official report and normalizing notation, context, and question form. These are supplied registry assertions, not independent verification here.
\item Supplied report citation: \emph{Invariants in Low-Dimensional Topology and Knot Theory (2013)}, Oberwolfach Reports, DOI \texttt{10.4171/owr/2012/28}, pp.\ 1687--1758. \url{https://doi.org/10.4171/owr/2012/28}. The supplied background describes Oberwolfach Reports as documenting mathematical workshops and their open-problem discussions.
\item Supplied literature reference: Manturov, \emph{Free Knots and Parity}. \url{https://arxiv.org/abs/0912.5348}.
\item Supplied literature reference: original Oberwolfach report. \url{https://doi.org/10.4171/owr/2012/28}.
\item Literature assessment dated exactly \texttt{2026-08-21}: the supplied curation labels the questions open and states that parity for knots in arbitrary \(3\)-manifolds and applications to link homotopy remain open-ended research questions, and that the cited free-knot results do not resolve these extensions. This is supplied dated literature triage, not an independently verified assessment; the record requests independent verification of time-sensitive status and cited claims.
\end{itemize}
\end{document}
